\begin{proof}
\begin{enumerate}[label=\arabic*.]
\item \textbf{The actual height, parameter, and real coordinate maps.}
Put $k=625\pi/96$, $c_0=k^{-1}$, and use the specified right
connection $H'=H\Gamma_6$, $H(1)=I$, with $Y=XH$.
For the uniquely normalized real quotient-intersection vector write
\[
 q=(0,v_2,v_1,-i,0,0)^t,\qquad z=Yq,
 \qquad h_b=c_0\operatorname{Im}z_0.
\]
The lower five entries of $z$ are real. Every marked monodromy
matrix $T$ is real and fixes $e_0$. Its top entry on $z$ is
$z_0$ plus a real linear combination of those lower entries, so
$\operatorname{Im}(Tz)_0=\operatorname{Im}z_0$. This proves
branch independence in the given marking.
The normalization gives $\alpha(q)=k$, since the coefficient of
$w_0$ is $-i$ and $\alpha(w_0)=ki$. The lifted vector for parameter
$\xi$ therefore has top coordinate $z_0+k\xi$ and the same lower
coordinates. The exact obstruction in the real quotient-intersection
line is
\[
 \operatorname{Im}(z_0+k\xi)=k(h_b+\operatorname{Im}\xi).
\]
The transversality criterion of \cref{night:P1.8} consequently
uses this same function for complex and real parameters.

The numerical preconditioner has an explicit real coordinate map.
For a constant limiting root frame $H_c$, let $B$ have lower-five
columns
\[
 \operatorname{Im}(H_c)_1,\quad\operatorname{Re}(H_c)_1,
 \quad\operatorname{Im}(H_c)_2,\quad\operatorname{Re}(H_c)_2,
 \quad\operatorname{Im}(H_c)_3.
\]
The certified inverse of $B$ defines
\[
 r=(\operatorname{Im}(H_c)_{01},\operatorname{Re}(H_c)_{01},
       \operatorname{Im}(H_c)_{02},\operatorname{Re}(H_c)_{02},0)B^{-1},
 \qquad
 T_c=\begin{pmatrix}1&-r\\0&B^{-1}\end{pmatrix}.
\]
This is a real invertible matrix fixing $e_0$.
Realness of $B^{-1}z_{\rm low}$ is equivalent to realness of
$z_{\rm low}$, and the change of top entry is real on that vector.
Thus it preserves both the normalized pair and $h_b$ exactly.
Writing $v_2=a+ib$, $v_1=c+id$ gives the actual real equations
with row
$(\operatorname{Im}Y_{i1},\operatorname{Re}Y_{i1},
\operatorname{Im}Y_{i2},\operatorname{Re}Y_{i2})$
and right-hand entry $\operatorname{Re}Y_{i3}$.
Approximate row selection only chooses four rows; the inverse
and every ensuing enclosure are certified in ball arithmetic.

\item \textbf{The marked conjugation map.}
The connection has real rational coefficients. Conjugating the
specified ordered polygonal meridians gives their inverse paths.
The actual matrix $K_0=X\overline X^{-1}$ therefore satisfies
$K_0A=A^{-1}K_0$, $K_0B=B^{-1}K_0$, and $K_0e_0=e_0$.
The normalized intertwiner uniqueness in \cref{night:P1.15}
applies: the quotient of two invertible solutions is in the
common centralizer fixing $e_0$. Direct exact solution gives
\[
 K=\begin{pmatrix}
 1&2&-1&-1&2&-1\\
 0&0&0&0&-1&1\\
 0&0&-1&0&0&1\\
 0&0&0&-1&0&1\\
 0&-1&0&0&0&1\\
 0&0&0&0&0&1
 \end{pmatrix}.
\]
Multiplication verifies the two intertwiner equations,
$Ke_0=e_0$, $K^2=I$, and $\det K=-1$, hence $K=K_0$.
On conjugate continuation branches,
$Y(\overline\lambda)=K\overline{Y(\lambda)}$.
The vector $q'=-\overline q$ has the required $-i$ normalization,
and its image is $-K\overline z$, with real lower coordinates.
Uniqueness therefore identifies the normalized intersection
vector and gives
$h_b(\overline\lambda)=h_b(\lambda)$.
This is the actual coordinate map used in reflecting the covered
regions.

\item \textbf{Rigorous analytic enclosures.}
The following formulas specify the complete-domain computation.
All matrix coefficients come from the displayed rational
connection and the specified gauges of
\cref{night:P1.15,night:P1.16}. A regular root frame is the
pullback followed by its stated diagonal gauge $D$; its right
connection is exactly
$D^{-1}\lambda'\Gamma_6D+D^{-1}D'$.
At a positive anchor the constant limiting frame is obtained by
the actual map $H_c=Y_{\rm root}(z_c)U(z_c)^{-1}$, using the
normalized continuation and a certified matrix inverse.

\paragraph{The analytic enclosures and their actual matrix maps.}
For positive weights $w_i$ let $W=\operatorname{diag}(w_i)$. The
specified norm is the maximum absolute row-sum norm of the actual
conjugate $W^{-1}MW$:
\[
 \|M\|_w=\max_i\sum_j|M_{ij}|\frac{w_j}{w_i}.
\]
The equality $W^{-1}MNW=(W^{-1}MW)(W^{-1}NW)$ and the rowwise
triangle inequality prove
$\|MN\|_w\le\|M\|_w\|N\|_w$ and
$|M_{ij}|\le\|M\|_w w_i/w_j$. For matrices with centres $C,D$
and entrywise error bounds $E,F$, multiplication has entrywise
error at most
\[
 |C|F+E|D|+EF,
\]
with the computed rounding error also retained. In particular, a
left factor $H_c$ acts on an entrywise error matrix by the actual
sum $|H_c|_+E$, where $|H_c|_+$ bounds its entries including their
ball radii.

\paragraph{Root series.}
The listed regular-root connection matrices already include the
specified diagonal gauges. Their frame map is
$E_{\rm root}=E_{\rm original}D_{\rm root}$, and the numerical
coordinate change is exactly
$G_z(z)=\alpha G_{\rm root}(\alpha z)$, where
$\alpha^4=256/3125$ at zero and $\alpha^5=3125/256$ at infinity.
Multiplying $G_{\rm root}(t)$ by $1+t^m/\alpha^m$ gives the exact
rational polynomial $\sum_{j=0}^{m-1}C_jt^j$, so
$B_j=\alpha^{j+1}C_j$. Thus
\[
 U'=U\frac{\sum_{j=0}^{m-1}B_jz^j}{1+z^m},\qquad U(0)=I.\tag{1.18.2}
\]
Coefficient comparison in $(1+z^m)U'=UB$ gives
\[
 (n+1)U_{n+1}
 =\sum_{j=0}^{\min(m-1,n)}U_{n-j}B_j
   -(n-m+1)U_{n-m+1},
\]
where the last term is absent for $n<m$. The order of every
matrix product is as displayed. Picard iteration on any smaller
disc converges by the bound $(Lr)^k/k!$ on its $k$-fold integral,
so these are the coefficients of the actual normalized solution.

Use $w=(1,1/128,1/128,1/32,1/64,1/32)$ and
\[
 \mathfrak b=\max_i\sum_{j,k}|(B_k)_{ij}|\frac{w_j}{w_i}.
\]
On a ray of radius $t<1$ the connection norm is at most
$\mathfrak b/(1-t^m)$. The scalar Gronwall inequality
\cite[Lemma~2.7, pp.~42--43]{TeschlODE2012} applies to the actual
matrix integral equation. The elementary estimate
\[
 \frac1{1-t^m}\le\frac1{m(1-t)}+\frac{m-1}{m}
\]
holds because its numerator difference is
$\bigl(\sum_{j=1}^{m-1}t^j-(m-1)t^m\bigr)/m\ge0$.
Consequently
\[
 \|U(z)\|_w\le K(R):=
 \exp\!\left\{\mathfrak b\left[-\frac{\log(1-R)}m
                 +\frac{m-1}{m}R\right]\right\}
 \quad (|z|\le R<1).
\]
Cauchy's coefficient integral, with the positively oriented circle
of radius $R$ \cite[\S1.9(iii), (1.9.30)--(1.9.31)]{DLMFCauchy2026},
gives $\|U_n\|_w\le K(R)R^{-n}$. Hence the error after degree $N$
on the whole disc $|z|\le r<R$ is
\[
 \tau_N(r;R)=K(R)\frac{(r/R)^{N+1}}{1-r/R}.\tag{1.18.3}
\]
Its $ij$ entry is bounded by $\tau_Nw_i/w_j$; after left
multiplication by $H_c$ the bound is
$\tau_N\sum_l|(H_c)_{il}|_+w_l/w_j$.
The current root computations use $N=3000$, $R=0.999$ and each
rectangle's actual maximum radius. Every coefficient and point
rounding radius remains in the corresponding ball calculation.

\paragraph{Enclosing an entire rectangle.}
Let $z_c$ be its centre, $\rho$ its Euclidean half-diagonal, and
$\mathcal A$ a nonnegative entrywise bound for $|G(z)|$ on the
whole rectangle. For any point $z$ the segment
$z_c+t(z-z_c)$ remains in it. Iteration of the integral equation
on that segment bounds the $k$-th ordered integral by
$|Y(z_c)|_+|z-z_c|^k\mathcal A^k/k!$. Subtracting its zeroth term
and summing proves the entrywise simultaneous enclosure
\[
 |Y(z)-Y(z_c)|
 \le |Y(z_c)|_+\bigl(\exp(\rho\mathcal A)-I\bigr).
\]
The centre enclosure and matrix-exponential rounding radii are
retained, so this bounds every point of the rectangle.

The midpoint-polynomial bound has the same explicit domain
control. For $p(z)=\sum_{j=0}^Nc_jz^j$, coefficient discs with
midpoints $\widehat c_j$ and radii $\epsilon_j$, and midpoint
evaluation rounding radius $\eta$, the centre error is
$\eta+\sum_j\epsilon_j|z_c|^j$. Integrating $p'$ on each segment
bounds the remaining variation by
\[
 \rho\sum_{j=1}^N j\bigl(|\widehat c_j|+\epsilon_j\bigr)r^{j-1},
\]
where $r$ is the rectangle's maximum modulus. Add the analytic
tail at that same $r$. This formula includes coefficient
uncertainty in the derivative bound as well as at the centre.

\paragraph{The cusp holomorphic factor.}
The exact connection here is
$R/q+A_0/(1+q)$, with $R^2=0$. In
$Y=H(I+R\log q)P(q)$, $P(0)=I$, the inverse of the logarithmic
factor is $I-R\log q$. Multiplication by that inverse gives
\[
 P'=\frac{\operatorname{ad}(P)}q+\frac{PA_0}{1+q},
 \qquad\operatorname{ad}(Q)=QR-RQ.
\]
Thus
\[
 (nI-\operatorname{ad})P_n
  =\sum_{j=0}^{n-1}(-1)^{n-1-j}P_jA_0.
\]
Equivalently its right side is
$\operatorname{ad}(P_{n-1})+P_{n-1}A_0-(n-1)P_{n-1}$.
Direct expansion gives
$\operatorname{ad}^2(Q)=-2RQR$ and
$\operatorname{ad}^3(Q)=0$, so the actual inverse is
\[
 (nI-\operatorname{ad})^{-1}
  =\frac1n\left(I+\frac{\operatorname{ad}}n
                   +\frac{\operatorname{ad}^2}{n^2}\right).
\]
This fixes both the commutator sign and the side of the $A_0$
multiplication.

Use the actual cusp weights $(100,1,1,1,1,1)$ and put
$r_0=\|R\|_w$, $a_0=\|A_0\|_w$,
$C_0=1+2r_0+4r_0^2$, $\gamma=C_0a_0$. For the supplied exact
matrices,
\[
 r_0=\frac{5103}{8000},\quad a_0=\frac{367}{512},\quad
 C_0=\frac{62452609}{16000000},\quad
 \gamma=\frac{22920107503}{8192000000}.
\]
The norm of the inverse operator is at most $C_0/n$. Therefore
$\|P_n\|_w\le(\gamma/n)\sum_{j<n}\|P_j\|_w$.
The nonnegative scalar majorant $c_0=1$,
$c_n=(\gamma/n)\sum_{j<n}c_j$ satisfies
$c_n=(\gamma+n-1)c_{n-1}/n$ and has generating function
$(1-s)^{-\gamma}$. Its radius is at least one by the coefficient
ratio. Thus the actual matrix series converges normally on
$|q|<1$, solves the equation term by term, and obeys
\[
 \|P(q)\|_w\le K_0(s):=(1-s)^{-\gamma}\quad(|q|\le s<1).
\]
The coefficient integral and geometric sum give its degree-$N$
tail on $|q|\le r<s$ as
\[
 K_0(s)\frac{(r/s)^{N+1}}{1-r/s},
\]
with the exact entry factor $w_i/w_j$. The current cusp series
uses $N=600$ and $s=0.4$ for $r\le0.25$. On the inner disc
$r=0.01$ the enclosure retains all actual coefficient absolute
sums through degree600 and this tail. It bounds $P$ itself;
the subsequent logarithmic factor retains its specified branch
and its own actual matrix entries.

\paragraph{Collar continuation.}
At a centre $z_c$ put $t=z-z_c$ and expand
\[
 1+(z_c+t)^5=\sum_{j=0}^5q_jt^j,
 \quad q_j=\binom5jz_c^{5-j}+\delta_{j0},
 \qquad
 \widetilde B_j=\sum_{k=j}^4\binom kjB_kz_c^{k-j}.
\]
For $T(0)=I$ the exact transition coefficients satisfy
\[
 T_{n+1}=
 \frac{\displaystyle
   \sum_{j=0}^{\min(4,n)}T_{n-j}\widetilde B_j
  -\sum_{j=1}^{\min(5,n)}q_j(n-j+1)T_{n-j+1}}
 {(n+1)q_0}.
\]
Let $\zeta_j=e^{(2j+1)\pi i/5}$ and
$R_{\rm loc}=\tfrac12\min_j|z_c-\zeta_j|$. On the closed disc
of this radius,
\[
 |1+(z_c+t)^5|\ge D_{\rm lower}
   :=\prod_j\bigl(|z_c-\zeta_j|-R_{\rm loc}\bigr)>0.
\]
The numerator norm is at most
\[
 B_{\max}=\max_i\sum_{j,k}|(\widetilde B_k)_{ij}|
                  \frac{w_j}{w_i}R_{\rm loc}^k.
\]
The actual complete-domain code uses the further upper bound
\[
 \widehat B_{\max}=\sum_{k=0}^4\|B_k\|_w
                           (|z_c|+R_{\rm loc})^k.
\]
To see the exact comparison, substitute the displayed formula for
$\widetilde B_l$ in each row sum and use the triangle inequality.
Reindex the finite set $0\le l\le k\le4$ by $(l,k)\mapsto(k,l)$.
The sum in $l$ is
$\sum_{l=0}^k\binom kl|z_c|^{k-l}R_{\rm loc}^l
=(|z_c|+R_{\rm loc})^k$.
Each remaining row sum is at most $\|B_k\|_w$; taking the row
maximum afterward proves $B_{\max}\le\widehat B_{\max}$.
Thus the bound actually used below is
$M=\widehat B_{\max}/D_{\rm lower}$. Radial Gronwall and Cauchy
give the exact degree-$N$ transition tail
\[
 e^{MR_{\rm loc}}
 \frac{(|\delta|/R_{\rm loc})^{N+1}}
      {1-|\delta|/R_{\rm loc}}.
\]
The current steps use $N=160$ and
$|\delta|/R_{\rm loc}<0.3$. The actual continuation order is
$Y(z_c+\delta)=Y(z_c)T_{z_c}(\delta)$, as follows by checking
the right ODE and its initial value. Every transition tail is
added entrywise before multiplying on the right, and the
entrywise product rule propagates all earlier errors.

\paragraph{The real implicit system and Hessian.}
Use the specified fixed four rows on the central square and
write the actual system as $Mv=b$, with
\[
 M_{i,*}=(\operatorname{Im}Y_{i1},\operatorname{Re}Y_{i1},
          \operatorname{Im}Y_{i2},\operatorname{Re}Y_{i2}),
 \qquad b_i=\operatorname{Re}Y_{i3}.
\]
For the top row define $c$ by the same four coordinate projections,
$d=\operatorname{Re}Y_{03}$ and $h=c_0(cv-d)$,
$c_0=96/(625\pi)$. The certified determinant enclosure excludes
zero on the entire central square and on its sixteen subsquares;
the matrix inverse is therefore the actual one there.
For $a,b\in\{x,y\}$, direct differentiation gives
\[
 \begin{split}
 v_a&=M^{-1}(b_a-M_av),\\
 v_{ab}&=M^{-1}(b_{ab}-M_{ab}v-M_av_b-M_bv_a),\\
 h_a&=c_0(c_av+cv_a-d_a),\\
 h_{ab}&=c_0(c_{ab}v+c_av_b+c_bv_a+cv_{ab}-d_{ab}).
 \end{split}
\]
The repeated-index cases retain both cross terms. The actual
complex derivatives supplying these real projections are
\[
 Y_x=YG,\quad Y_y=iYG,\quad
 Y_{xx}=Y(G^2+G'),\quad Y_{xy}=iY(G^2+G'),\quad
 Y_{yy}=-Y(G^2+G').
\]
Taking the stated real and imaginary parts and carrying out these
products with rigorous full-square balls consequently encloses
all Hessian entries throughout the domain. Approximate row
selection merely chooses a minor; it supplies neither its inverse
nor its error bound.

\item \textbf{The infinity value and its strict local minimum.}
Put $a=256/3125$, $\lambda=u^{-5}$ and
\[
 D_\infty(u)=\operatorname{diag}(1,u^{-8},u^{-4},1,u^4,u^8),
 \qquad \widetilde Y(u)=XH(u^{-5})D_\infty(u).
\]
The specified root extension makes $\widetilde Y$ holomorphic
and invertible at zero. The $w_0$ column normalization is unchanged.
In the preconditioned root frame the chosen four-by-four real
system has determinant in
\[
 [\,0.997929591387219260,\ 1.002070408612780740\,]
\]
on the entire central square, with its inverse certified there.
Its inverse defines a real analytic normalized coefficient pair
at zero. The fifth reality equation extends by continuity from
the punctured chart, and any solution must satisfy the four
invertible equations. This proves existence and uniqueness there.

For $\zeta=e^{2\pi i/5}$ set
$T_\zeta=\operatorname{diag}(1,\zeta^{-8},\zeta^{-4},1,
\zeta^4,\zeta^8)$. The actual deck formula is
\[
 \widetilde Y(\zeta u)=A\widetilde Y(u)T_\zeta.
\]
At zero it gives
$A\widetilde Y(0)=\widetilde Y(0)T_\zeta^{-1}$.
Applying $A$ to the normalized real-lower vector changes its
coefficient pair to $(\zeta^8v_2,\zeta^4v_1)$ and preserves
the $-i$ coefficient. Uniqueness forces $v_2=v_1=0$.
Consequently
\[
 h_b(\infty)=-c_0\operatorname{Re}\widetilde Y(0)_{03}
            =-\frac{96}{625\pi}\operatorname{Re}(XH_\infty)_{03}.
\]
For general $u$, the same deck formula and the real monodromy
map give $h_b(\zeta u)=h_b(u)$. Its real differential at zero
is fixed by a rotation through $2\pi/5$, so it vanishes.

The numerical infinity coordinate is exactly
$w=a^{1/5}u$, hence $\lambda=aw^{-5}$.
The sixteen closed subsquares cover
$Q=[-1/10,1/10]^2$ in $(\operatorname{Re}w,\operatorname{Im}w)$.
The implicit derivative formulas above give on their union
\[
 h_{xx}>0.3215,\qquad h_{yy}>0.3220,
 \qquad |h_{xy}|<0.0218.
 \tag{1.18.5}
\]
For every real vector $t=(t_1,t_2)$, the inequality
$2|t_1t_2|\le t_1^2+t_2^2$ yields
$t^tD^2h_b\,t>(299/1000)|t|^2$.
The gradient is still zero after the positive scalar coordinate
change. Integrating twice along the segment in the convex square
therefore gives
\[
 h_b(aw^{-5})-h_b(\infty)
 =\int_0^1(1-s)w^tD^2h_b(sw)w\,ds
 \ge\frac{299}{2000}|w|^2.
\]
Thus the center is the unique minimum on $Q$.

The independent normalized continuations give the following
outward enclosure, whose endpoints are exact terminating decimals:
\[
 \begin{gathered}
 -1.094540250311402780116268263054945703726
       <h_b(\infty),\\
 h_b(\infty)<
 -1.094540250311402780116268263054945703725.
 \end{gathered}
 \tag{1.18.4}
\]
The complete stored enclosure has radius less than
$4.38\cdot10^{-100}$. It lies in the statement's earlier
$4.04\cdot10^{-56}$ enclosure. The independently authored
root-connection computation agrees with the independently
propagated normalized period calculation to more than89 digits;
both enclose the exact expression just proved.

\item \textbf{The closed covers and their strict lower bounds.}
Let $R=49/50$. The infinity sector
$0\le\arg w\le2\pi/5$, $|w|\le R$ maps under
$\lambda=aw^{-5}$ onto $|\lambda|\ge aR^{-5}$.
The square $[0,1/10]^2$ removed from its leaf calculation lies
in $Q$ and has already been treated.
At zero put $t=a^{1/4}v$, so $\lambda=t^4=av^4$.
The first-quadrant sector $|v|\le R$ covers
$|\lambda|\le aR^4$, including the completed zero root.
The marked real matrices $A$ and $B$ identify the respective
sector boundaries and preserve the height.

The intervening radial region is precisely the image of
\[
 R\le|w|\le R^{-4/5},\qquad 0\le\arg w\le2\pi/5.
\]
Indeed its endpoint radii map to $aR^{-5}$ and $aR^4$.
Here the cusp coordinate is exactly
$q=\lambda/\lambda_*-1=-w^{-5}-1$.
The collar certificate treats its closed part $|q|\ge1/4$.
The remaining part is in the cusp disc; the annulus certificate
treats $1/100\le|q|\le1/4$ for all arguments. The punctured
inner disc is treated in the next item. These closed regions
include their shared boundaries and cover the whole punctured
base, with the two finite-root points included as specified.

For each required rectangle, a whole-domain enclosure is either
strictly accepted or the rectangle is subdivided. A discarded
rectangle has empty intersection with the required region or is
contained in the already treated central square. The complete
finite subdivision and domain checks have no unresolved leaves.
The current computation uses110-digit ball arithmetic. Its
unrounded lower endpoints have been exported and replayed against
the same exact rational threshold
\[
 U=-1.0945402503114027801162682630549457037
       >h_b(\infty).
\]
The following terminating decimals are exact uniform lower
bounds, strictly smaller than every recorded lower endpoint
in the indicated region and strictly greater than $U$:
\[
 \begin{array}{lrr}
 \text{region}&\text{accepted leaves}&\text{strict lower bound for }h_b\\
 \hline
 \text{infinity sector outside the square}&1203&-1.094400523035\\
 \text{zero sector}&592&-1.083329522544\\
 \text{cusp annulus}&10780&-1.094402439774\\
 \text{collar outside the cusp disc}&336&-1.092202452969
 \end{array}
\]
All12911 precise endpoint comparisons pass. The endpoint replay
uses the same newly authored complete-domain computation and
retains its full interval endpoints. It adds a directly
replayable inequality record to that computation.

\item \textbf{The actual cusp graph and the inner disc.}
Use $\lambda=\lambda_*(1+q)$ and the actual diagonal gauge
\[
 G=Y\operatorname{diag}(1,\lambda^2,\lambda,1,
                         \lambda^{-1},\lambda^{-2})
   =H_c(I+R_c\log q)P(q).
\]
The cusp series and continuation above give the normalized
$H_c$ and $P$. The matrix $N=2\pi iH_cR_cH_c^{-1}$ is the
logarithm of the actual marked cusp monodromy. Its certified
enclosure contains the following unique integer matrix, and
exact multiplication verifies the stated marking:
\[
 N=\begin{pmatrix}
 0&0&0&-1&1&-1\\0&0&0&0&0&0\\0&0&0&0&0&0\\
 0&1&0&0&0&0\\0&0&1&0&0&0\\0&-1&1&0&0&0
 \end{pmatrix},\qquad
 J=\begin{pmatrix}
 -1&0&0&0&0&0\\0&0&0&0&1&0\\0&0&0&0&0&1\\
 0&1&0&1&0&0\\0&0&1&0&0&0\\0&-1&1&0&0&0
 \end{pmatrix}.
\]
In particular $N^2=0$, $J$ is unimodular,
$J^{-1}NJ=\left(\begin{smallmatrix}0&I_3\\0&0\end{smallmatrix}\right)$,
and $e_J:=J^{-1}e_0=(-1,0,0,0,0,0)^t$.
Put $c_*=-h_b(\infty)$, $s=kc_*$, and $V=J^{-1}H_cP(q)$.
The three actual period columns are
$V_1,V_2,V_3-se_J$. Their upper and lower blocks will be
denoted by $U_0,L$. The right scaling from the algebraic period
columns is exactly $\operatorname{diag}(\lambda^2,\lambda,1)$;
it preserves their span and the third-column shift.
The inverse of $L$ is certified on the entire closed disc
$|q|\le1/100$.
Since $Ne_0=0$, the logarithmic factor acts on these shifted
columns by
$\left(\begin{smallmatrix}I&\ell I\\0&I\end{smallmatrix}\right)$,
where $\ell=\log q/(2\pi i)$. The actual graph is therefore
\[
 Z(q)=\frac{\log q}{2\pi i}I_3+C(q),\qquad C(q)=U_0L^{-1}.
\]
For $B(q)=\operatorname{Im}C(q)$ and
$y=-\log|q|/(2\pi)$, expansion of the determinant gives
\[
 \det\operatorname{Im}Z
 =y^3+\operatorname{tr}(B)y^2
  +\sum_{i<j}(B_{ii}B_{jj}-B_{ij}B_{ji})y+\det B.
 \tag{1.18.6}
\]
Put $y_0=\log100/(2\pi)$ and $t=y-y_0\ge0$.
The four coefficients of
$\det(tI_3+y_0I_3+B(q))$, in increasing degree, have the
uniform lower bounds
\[
 1.9918,\qquad4.9602,\qquad3.9580,\qquad1,
\]
with the last coefficient exactly one. Thus the determinant is
positive on the punctured inner disc.
In the actual real $J$ coordinates, a graph vector $(Zv,v)$
is real precisely when $v$ is real and $(\operatorname{Im}Z)v=0$.
The determinant therefore proves transversality of the
$c_*$ plane there. By the first item's exact obstruction map,
$h_b-h_b(\infty)$ never vanishes. The punctured disc is connected,
and at $q=1/100$ the certified value
$h_b=-0.4504422009074\ldots$ is greater than $h_b(\infty)$.
Consequently $h_b>h_b(\infty)$ throughout this disc.

For completeness these same graph maps prove the cusp limit.
The lower block of $e_J$ vanishes, so $L$ is independent of a
real parameter $c$, and $C_c$ is holomorphic at zero and affine
in $c$. For any fixed real bound $M$, it is uniformly bounded
on a smaller closed disc and the compact interval between
$c_*$ and $-M$. As $y\to\infty$, the Neumann inverse makes
$yI+\operatorname{Im}C_c$ invertible uniformly on that interval.
The real quantity $h_b+c$ cannot cross zero along it. Its sign
at $c_*$ is positive, so $h_b-M>0$ at the other endpoint.
This proves $h_b\to+\infty$ at the cusp. The covers and the
strict local estimate now prove the global minimum and its
uniqueness.

\item \textbf{The parameter domain and the full splitting translation.}
The first item's obstruction map and \cref{night:P1.8} give
$c>-h_b(\infty)$ and $\operatorname{Im}\xi>-h_b(\infty)$
on the completed base. The enclosure above gives
$-h_b(\infty)<1.095<10$.
For the stated upper-half-plane polygon, let
$g=(96/(625\pi i))w_0$, so $\alpha(g)=1$.
The exact top coordinate of its algebraic lift is
\[
 a_{\rm split}=\frac{96}{625\pi i}Y_{03}(\lambda_*/2).
\]
The defining equality
$\sigma_{\rm original}=\sigma_b-a_{\rm split}\alpha e$
gives the actual plane equality
\[
 (\sigma_{\rm original}+\xi_{\rm original}\alpha e)(P)
  =(\sigma_b+(\xi_{\rm original}-a_{\rm split})\alpha e)(P).
\]
Hence $\xi_b=\xi_{\rm original}-a_{\rm split}$ and, on the
normalized vector with $\alpha(q)=k$,
$h_{\rm original}=h_b-\operatorname{Im}a_{\rm split}$ pointwise.
The continuation gives the stated real and imaginary parts of
$a_{\rm split}$ with radii less than $4.41\cdot10^{-100}$ and
$1.41\cdot10^{-100}$, respectively. The resulting minimum is
\[
 -0.181138544374987443901945340672108376803615204418668779889\ldots,
\]
with stored radius less than $5.77\cdot10^{-100}$, attained
at the same base point. The original real line
$\xi_{\rm original}=ic$ maps exactly to
\[
 \xi_b=-\operatorname{Re}a_{\rm split}
          +i(c-\operatorname{Im}a_{\rm split}).
\]
Both coordinates are retained, and its criterion is precisely
$c>-h_{{\rm original},\min}$.

The exact coefficient arrays, finite subdivision trees, actual
matrix inverses and unrounded endpoint certificates are retained
in \path{code/continuation_V/}; in particular the complete
region summary is
\path{minimum_precise_endpoint_summary.json} there.
The separate analytic certificate uses256-bit balls and713
checks of the displayed recurrences and matrix maps; its tested
representative domains are recorded separately from the complete
110-digit leaf partition used above.
\end{enumerate}
\end{proof}
